En este trabajo se estudió la síntesis y estructura cristalina de nanoplatos triangulares de plata (AgTNPs). Las partículas fueron sintetizadas mediante reducción química en solución acuosa y caracterizadas por microscopía electrónica de transmisión (TEM). Las imágenes de campo claro muestran una morfología triangular con lados entre 20 y 120~nm y una
orientación preferencial con las caras $\{111\}$ paralelas al sustrato. El análisis por transformada de Fourier de imágenes de alta resolución (HRTEM) reveló reflexiones $\frac{1}{3}\hkl{422}\}$, características de regiones con estructura hcp originadas por fallas de apilamiento en los planos $\{111\}$, con un espaciado interplanar de \SI{0.26 \pm 0.1}{\nano\metre} consistente con los valores reportados en la literatura. Para estudiar la relación entre las fallas y el contraste en imágenes de HRTEM se realizaron simulaciones por el método multislice . Los resultados muestran que una estructura fcc pura resulta prácticamente invisible en esta condición de imagen, lo que implica que las partículas observadas contienen necesariamente fallas de apilamiento. El estudio de estructuras tipo ``sánguche'' evidenció además que el contraste final depende de manera compleja de la distribución global de las fallas, yno únicamente de su cantidad, indicando la necesidad de modelos estructurales más completos para la interpretación de las imágenes experimentales.